\documentclass[portrait,final,a0paper]{baposter}
%\documentclass[a4shrink,portrait,final]{baposter}
% Usa a4shrink for an a4 sized paper.

\tracingstats=2
\usepackage{calc}
\usepackage{graphicx}
\usepackage{amsmath}
\usepackage{amssymb}
\usepackage{relsize}
\usepackage{multirow}
\usepackage{bm}

\usepackage{graphicx}


\usepackage{pgfbaselayers}
\pgfdeclarelayer{background}
\pgfdeclarelayer{foreground}
\pgfsetlayers{background,main,foreground}

\usepackage{times}
\usepackage{helvet}
%\usepackage{bookman}
\usepackage{palatino}
\usepackage{background}

\newcommand{\captionfont}{\footnotesize}

\selectcolormodel{cmyk}

\graphicspath{{images/}}


%%%%%%%%%%%%%%%%%%%%%%%%%%%%%%%%%%%%%%%%%%%%%%%%%%%%%%%%%%%%%%%%%%%%%%%%%%%%%%
%%% Begin of Document
%%%%%%%%%%%%%%%%%%%%%%%%%%%%%%%%%%%%%%%%%%%%%%%%%%%%%%%%%%%%%%%%%%%%%%%%%%%%%%


\begin{document}
\typeout{Poster rendering started}	
\tikzset{ 
	table/.style={
		matrix of nodes,
		row sep=-\pgflinewidth,
		column sep=-\pgflinewidth,
		nodes={
			rectangle,
			draw=black,
			align=center
		},
		minimum height=1.5em,
		text depth=0.5ex,
		text height=2ex,
		nodes in empty cells,
		%%
		every even row/.style={
			nodes={fill=gray!20}
		},
		column 1/.style={
			nodes={text width=2em,font=\bfseries}
		},
		row 1/.style={
			nodes={
				fill=black,
				text=white,
				font=\bfseries
			}
		}
	}
}


\SetBgContents{}
%%% Setting Background Image %%%%%%%%%%%%%%%%%%%%%%%%%%%%%%%%%%%%%%%%%%%%%%%%%%
\background{
	\begin{tikzpicture}[remember picture,overlay]%
	\draw (current page.center)+(-1em,1em) node[anchor=center,opacity=0.2, scale=1.1]
	{\includegraphics[width=1.1\textwidth]{background1}};
	\end{tikzpicture}
}	

%%%%%%%%%%%%%%%%%%%%%%%%%%%%%%%%%%%%%%%%%%%%%%%%%%%%%%%%%%%%%%%%%%%%%%%%%%%%%%
%%% Here starts the poster
%%%---------------------------------------------------------------------------
%%% Format it to your taste with the options
%%%%%%%%%%%%%%%%%%%%%%%%%%%%%%%%%%%%%%%%%%%%%%%%%%%%%%%%%%%%%%%%%%%%%%%%%%%%%%
% Define some colors
\definecolor{DarkMagenta}{cmyk}{0,0.64,0.2,0}
\definecolor{yellow}{cmyk}{0.3,0.1,0.2,0.0}
\definecolor{reddishyellow}{cmyk}{0,0.22,1.0,0.0}
\definecolor{black}{cmyk}{0,0,0.0,1.0}
\definecolor{darkYellow}{cmyk}{0,0,1.0,0.5}
\definecolor{Carnelian}{cmyk}{1.0,0.29,0.0,0.26}

\definecolor{lightyellow}{cmyk}{0.0,0.0,0.0,0.0}
\definecolor{lighteryellow}{cmyk}{0,0,0.0,0.0}
\definecolor{lighteryellow}{cmyk}{0,0,0.0,0.0}
\definecolor{lightestyellow}{cmyk}{0,0,0.0,0.0}

\definecolor{primary1}{cmyk}{0.5,1.0,0.0,0.0}
\definecolor{primary2}{cmyk}{0,0,0,0.5}
\definecolor{primary3}{cmyk}{0,0,0,0.8}

\definecolor{secondary1}{cmyk}{0.82,0.38,0.2,0.2}
\definecolor{secondary2}{cmyk}{0.91,0.91,0.6,0.18}
\definecolor{secondary3}{cmyk}{0.71,0.23,0.13,0.3}
\definecolor{secondary4}{cmyk}{0.38,0.16,0.0,0.0}
\definecolor{secondary5}{cmyk}{0.20,0.62,0.2,0.0}
\definecolor{secondary6}{cmyk}{0.63,0.2,0.26,0.0}

%%
%\typeout{Poster Starts}

\newlength{\leftimgwidth}
\begin{poster}%
  % Poster Options
  {
  % Show grid to help with alignment
  grid=false,
  % Column spacing
  colspacing=1em,
  % Color style
  bgColorOne=lightestyellow,
  bgColorTwo=lightestyellow,
  borderColor=primary2,
  headerColorOne=primary1,
  headerColorTwo=reddishyellow,
  headerFontColor=white,
  boxColorOne=lightyellow,
  boxColorTwo=lighteryellow,
  % Format of textbox
  textborder=roundedleft,
%  textborder=rectangle,
  % Format of text header
  eyecatcher=false,
  headerborder=open,
  headerheight=0.095\textheight,
  headershape=roundedright,
  headershade=plain,
  headerfont=\Large\sf\bf, %Sans Serif
 % boxshade=plain,
%  background=shade-tb,
  background=user,
  linewidth=1.6pt
  }
  % Eye Catcher
  {\includegraphics[width=10em]{background1.png}} % No eye catcher for this poster. (eyecatcher=no above). If an eye catcher is present, the title is centered between eye-catcher and logo.
  % Title
  {
  \vspace{2.2em}
  \color{primary1}
  \textbf{
  Classification of Cells Based on Geometrical Features and Supervised Machine Learning}
  }
  % Authors
  {
  % Serif
  \vspace{0.5em}\textbf{\ S. Vohra\textsuperscript{2} and D. Prodanov\textsuperscript{1}}\\
   \textsuperscript{1}EHS, Imec, Leuven, Belgium, \textsuperscript{2} DTAI, KU Leuven\\
   \vspace{1.94em}
    
  }
  % University logo
  {% The makebox allows the title to flow into the logo, this is a hack because of the L shaped logo.  
  	 \makebox[11em][r]{%
  	 	\begin{minipage}{16em}
  	 		\hfill
  	 		\includegraphics[height=7.8em]{imec_logo}
  	 	\end{minipage}
  	 }
  }
   
    
%%%%%%%%%%%%%%%%%%%%%%%%%%%%%%%%%%%%%%%%%%%%%%%%%%%%%%%%%%%%%%%%%%%%%%%%%%%%%%
%%% Now define the boxes that make up the poster
%%%---------------------------------------------------------------------------
%%% Each box has a name and can be placed absolutely or relatively.
%%% The only inconvenience is that you can only specify a relative position 
%%% towards an already declared box. So if you have a box attached to the 
%%% bottom, one to the top and a third one which should be in between, you 
%%% have to specify the top and bottom boxes before you specify the middle 
%%% box.
%%%%%%%%%%%%%%%%%%%%%%%%%%%%%%%%%%%%%%%%%%%%%%%%%%%%%%%%%%%%%%%%%%%%%%%%%%%%%%
    %
 
%%%%%%%%%%%%%%%%%%%%%%%%%%%%%%%%%%%%%%%%%%%%%%%%%%%%%%%%%%%%%%%%%%%%%%%%%%%%%%
  \headerbox{Introduction}{name=contribution,column=0,row=0, span=1}{
%%%%%%%%%%%%%%%%%%%%%%%%%%%%%%%%%%%%%%%%%%%%%%%%%%%%%%%%%%%%%%%%%%%%%%%%%%%%%%
	\begin{itemize}
		\item Sparsity in Dataset.
		\item No full ground truth.	
		\item Subjectivity in expert assesments.
	\end{itemize}
		\includegraphics[width=7.2cm, height=2.8cm]{MONTAGE31}	
		\color{secondary1}\textbf{Anatomic \& TimeLapse Images for Segmentation }	
}

%%%%%%%%%%%%%%%%%%%%%%%%%%%%%%%%%%%%%%%%%%%%%%%%%%%%%%%%%%%%%%%%%%%%%%%%%%%%%%
\headerbox{Active Learning }{name=baseline,column=1,row=0,span=2}{
	\color{black}\text{Interactively collect new training examples by querying human user}
	\newline
	\includegraphics[width=15cm,height=5.3cm]{activelearning2}						
}


%%%%%%%%%%%%%%%%%%%%%%%%%%%%%%%%%%%%%%%%%%%%%%%%%%%%%%%%%%%%%%%%%%%%%%%%%%%%%%
  \headerbox{Features Extraction }{name=feature,column=0,span=2, below=contribution }{
%%%%%%%%%%%%%%%%%%%%%%%%%%%%%%%%%%%%%%%%%%%%%%%%%%%%%%%%%%%%%%%%%%%%%%%%%%%%%%
	\color{Carnelian}\textbf{FILTERS} \hspace{2.9cm} \textbf{CONVOLVED ROI} \hspace{1.9cm} \textbf{MULTIDIMENSIONAL FEATURE}
	\newline
	\color{Carnelian}\textbf{}\hspace{4.5cm}\textbf{AT VARIOUS SCALES} \hspace{1.9cm} \textbf{VECTOR OF EACH PIXEL}
	\newline 
		\begin{minipage}{.15\textwidth}	
			\color{black}\textbf{Gaussian }
			\newline
			\newline	
			\textbf{Gaussian X,Y}
			\newline 	
			\textbf{Derivatives} 
			\newline 
			\newline 
			\newline 
			\textbf{Laplace of} 
			\newline 
			\textbf{Gaussian}
			\newline  
			\newline
			\newline
			\textbf{BiLaplace of} 	
			\newline 
			\textbf{Gaussian}
		\end{minipage}
		\begin{minipage}{.15\textwidth}  
			\includegraphics[width=2.5cm, height=1.6cm]{gaussian}
			\includegraphics[width=2.5cm, height=1.6cm]{gaussiand}
			\includegraphics[width=2.5cm, height=1.6cm]{LOG}
			\includegraphics[width=2.5cm, height=1.6cm]{BLOG}	
		\end{minipage}
		\begin{minipage}{.10\textwidth}
			\includegraphics[width=1.3cm, height=1.2cm]{convolution1}
			\includegraphics[width=1.3cm, height=1.2cm]{convolution2}
			\newline 
			\color{primary1}\textbf{Region}
			\newline 
			\color{primary1} \hspace{0.3cm}\textbf{Of}
			\newline
			\textbf{Interest}
			\includegraphics[width=1.5cm, height=1.2cm]{convolution3}
			\includegraphics[width=1.5cm, height=1.2cm]{convolution4}
		\end{minipage}
        \begin{minipage}{.28\textwidth}
        	\includegraphics[width=4.1cm, height=5.8cm]{filteredImage}	
        \end{minipage}
        \begin{minipage}{.30\textwidth}
        	\includegraphics[width=4.2cm, height=2.9cm]{vector}	
        	\newline
        	\color{primary1} \hspace{1.7cm}\textbf{$\bigoplus$}
        	\newline
        	\color{Carnelian} \hspace{0.7cm}\textbf{CLASS LABEL}
        	\newline
        	\color{Carnelian} \hspace{0.3cm}\includegraphics[width=2.5cm, height=1.8cm]{pool}	
        \end{minipage}

  }



\headerbox{Results}{name=evaluation,column=0,span=3,below=feature}{
  \begin{minipage}{.20\textwidth}
  	\color{Carnelian}\hspace{0.1cm}\textbf{ACTIVE SEGMENTATION}
  		\begin{itemize}
  			\item ImageJ plugin.
  			\item Inspired by trainable Weka Segmentation.	
  		\end{itemize}	
  	\color{Carnelian}\hspace{1.5cm}\textbf{ASTROCYTES}
  	\newline
  	\newline
  	\includegraphics[width=4.5cm,height=4cm]{astrocyte}	
  	\newline
  	\color{Carnelian}\hspace{0.6cm}\textbf{FILTER SCREEN}
  	\newline
  	\newline
  	\includegraphics[width=5cm, height=4cm]{FilterScreen}
  	\newline	
  \end{minipage}
  	\begin{minipage}{.10\textwidth}
  		\includegraphics[width=2cm, height=2cm]{arrow1}
  	\end{minipage}
    \begin{minipage}{.28\textwidth}
    	\color{Carnelian}\hspace{0.6cm}\textbf{CONFIGURE SCREEN}
    	\newline
    	\newline
    	\includegraphics[width=6cm, height=5cm]{ClassScreen}
    	\newline
    	\newline
    	\color{Carnelian}\hspace{0.6cm}\textbf{TRAINING SCREEN}
    	\newline
    	\newline
    	\includegraphics[width=6cm, height=5cm]{FeatureScreen}
    		
    	
    \end{minipage}
     \begin{minipage}{.10\textwidth}
     	\includegraphics[width=2cm, height=2cm]{arrow1}
     \end{minipage}
    \begin{minipage}{.25\textwidth}	
	
    	\color{Carnelian}\hspace{0.6cm}\textbf{SEGMENTED ASTROCYTES}
    	\newline
    	\newline
    	\includegraphics[width=6cm, height=5cm]{astrocyteresult}
    	\newline	
    	\newline	
    	\color{Carnelian}\hspace{0.1cm}\textbf{EVALUATION}
    	\newline
    	Recall (R) is defined as the number of true positives ($T_p$) over the number of true positives plus the number of false negatives ($F_n$).
    	\begin{equation*}
    	\label{eq:scale}
    	{R} = \frac{T_p}{T_p + F_n}
    	\end{equation*}		
    \end{minipage}
	
}
  \headerbox{Conclusions \& Outlook}{name=conclusion,column=2,span=1,below=baseline}{
  	%%%%%%%%%%%%%%%%%%%%%%%%%%%%%%%%%%%%%%%%%%%%%%%%%%%%%%%%%%%%%%%%%%%%%%%%%%%%%%
  	Presented Approach is 
  	\begin{itemize}
  		\item \color{primary1}\textbf{Semi Supervised: }\color{black} The results can be refined with iterations.
  		\item \color{primary1}\textbf{Cell Type Independent: }\color{black} not only limited to specific type of cell.
  	\end{itemize} 
  \color{Carnelian}	Future efforts will include implementing the Structural SVM or Relational Learning for finding the suitable patterns in astrocytic arbors.
  }
	\hspace{5cm}
 \centerline{ 
 	\begin{minipage}{20em}
 		\includegraphics[height=3.5em]{google1}   		
 		\vspace{0.5cm}		
 	\end{minipage}
 	\begin{minipage}{20em}
 		\includegraphics[height=2.5em]{logo_kul}   		
 		\vspace{0.5cm}		
 	\end{minipage}
 		\begin{minipage}{20em}
 			\includegraphics[height=3.4em]{imec_logo}   		
 			\vspace{0.5cm}		
 		\end{minipage}
 }

\end{poster}

\end{document}
